%% notes-tex/030-scaling-study/sections/6-next.tex

\section{What to run next, by what it buys\secstatus{todo}}
\label{sec:next}

In order of the limit each one moves, with its cost in 48 h one-GPU jobs. None
is queued.

\pfstep{Read the calibration jobs before trusting the budget table} Two 3 h jobs
at 180x8 and 45x2 report records per second and the measured FLOPs per step at
both ends of the ladder. Equal paces mean the loader binds, every size costs
4.3 h an epoch, and the fixed-epoch slice of $L(N)$ is all 48 h can give; a
faster 45x2 means the small sizes can run to convergence on the full pool, and
the epoch budget should then scale per size by the ratio of paces. The
validation pass is also timed there: at 240 epochs of 1/64 it runs 240 times
over 37,705 triples at a validation batch of 32, and if one pass costs more than
a quarter of an epoch the small subsets need validation every second epoch to
finish inside 48 h, with the readout interpolating on the epochs that have a
validation row.

\pfstep{A 1/8 subset at 180x8, both builds} Two jobs. The subsets are nested by
one shuffle, so adding a denominator leaves the existing subsets unchanged and
adds a fifth point to the $D$ axis between 1/16 and 1/4, which is what the floored
form's held-out test needs (four points in sample) and what the gap needs for a
third converged point. Eighty epochs of 1/8 is 10.4M records, one 48 h job at the
assumed pace.

\pfstep{The learning rate at the top of the ladder} One or two jobs at 180x8 on
1/16 at 1e-4 (and 6e-4). The check at 45x2 on 1/4 says nothing about the end of
the ladder where the batch is four times smaller than the arm was tuned at, and
an $N$ exponent fit across sizes run at one learning rate is an exponent of the
schedule as much as of the capacity.

\pfstep{A converged full-data point} Not cheap: one hundred epochs at 180x8 is
nine chained 48 h jobs with a queue wait per link. Until it exists the top of
the $D$ axis and the whole $N$ axis are read at ten epochs, and the figures say
so. The alternative that costs nothing is to fit the converged curve on 1/64,
1/16, 1/8 and 1/4 only and treat the full point as the held-out prediction
target, which the readout already does when the full point is the largest $x$.

\pfstep{The token ablation} One 180x8 column of the 030 arm with
\texttt{dataset\_token.enabled=false} on the four subsets, four jobs, which
separates the token from the per-entry rows in the gap.

\pfstep{The compute axis} Three budgets spanning two decades of FLOPs, four
sizes at each, records seen set by the budget and the per-size cost per record
from the calibration jobs; twelve jobs, after the first two sweeps have shown
that the budget buys anything.
